\documentclass[11pt,a4paper]{article}
\usepackage[T1]{fontenc}
\usepackage{lmodern}
\usepackage[margin=26mm]{geometry}
\usepackage{amsmath,amssymb,amsthm,mathtools}
\usepackage{microtype}
\usepackage{booktabs,array}
\usepackage{xcolor}
\usepackage{enumitem}
\usepackage{fancyvrb}
\usepackage[colorlinks=true,linkcolor=blue!45!black,citecolor=blue!45!black,urlcolor=blue!45!black]{hyperref}
\numberwithin{equation}{section}
\newtheorem{theorem}{Theorem}[section]
\newtheorem{lemma}[theorem]{Lemma}
\newtheorem{proposition}[theorem]{Proposition}
\newtheorem{corollary}[theorem]{Corollary}
\theoremstyle{remark}
\newtheorem{remark}[theorem]{Remark}
\newcommand{\Z}{\mathbb Z}
\newcommand{\Q}{\mathbb Q}
\newcommand{\CT}{\operatorname{CT}}
\newcommand{\Diag}{\operatorname{Diag}}
\setlength{\emergencystretch}{2em}
\allowdisplaybreaks[2]
\title{Integral diagonal methods for\protect\\Sun's Conjecture 5.7(ii)}
\author{Research manuscript}
\date{7 November 2026}

\begin{document}
\maketitle
\begin{abstract}
Let
\(W_n(1,t)=\sum_{k=0}^{\lfloor n/2\rfloor}
t^k\binom n{2k}^2\binom{2k}k^2\).
We give complete proofs of the two assertions in Conjecture 5.7(ii)
of version 11 of Z.-W. Sun's 2010 preprint: for every positive integer
\(N\), the sum \(\sum_{n<N}(42n+37)(-1)^nW_n(1,-16)\) is divisible
by \(N\), and for every prime \(p\) it is congruent to
\(p(21(\frac p7)+16)\) modulo \(p^2\).
For the first assertion we extend an integral diagonal criterion to
the denominator \(1+u+v+\kappa uv\). At \(\kappa=2\), its odd
diagonals are controlled by Bing He's published divisibility theorem
for a Ramanujan-type cubic-binomial sum. Exact reflected products and
Abel summation establish every odd prime-power base. An integral
two-state digit operator propagates these bases to arbitrary odd
indices, while a termwise binomial divisibility handles every even
index. A separate finite proof establishes the prime-square
assertion, including \(p=2,3,7\). Its connection with established
results is made explicit. Historical priority is not asserted.
\end{abstract}

\section{The unchanged problem and the scope of the proof}

Write \(b_k=\binom{2k}{k}\), and define
\begin{equation}\label{eq:original-sequence}
 w_n=W_n(1,-16)=\sum_{k=0}^{\lfloor n/2\rfloor}
 (-16)^k\binom nk^2\binom{n-k}k^2
 =\sum_{k=0}^{\lfloor n/2\rfloor}(-16)^k b_k^2\binom n{2k}^2.
\end{equation}
The equality follows from
\(\binom nk\binom{n-k}k=\binom n{2k}\binom{2k}k\).
Set
\begin{equation}\label{eq:target-sum}
 S(N)=\sum_{n=0}^{N-1}(42n+37)(-1)^n w_n.
\end{equation}

\begin{theorem}\label{thm:main}
For every integer \(N\ge1\),
\begin{equation}\label{eq:main-allN}
                         N\mid S(N).
\end{equation}
For every prime \(p\),
\begin{equation}\label{eq:main-p2}
 S(p)\equiv p\left(21\left(\frac p7\right)+16\right)\pmod{p^2},
\end{equation}
where \((\frac p7)\) is the Legendre symbol, including the value
zero at \(p=7\).
\end{theorem}

These are exactly the two assertions of Sun's Conjecture 5.7(ii)
in \cite[Conjecture 5.7(ii), p.~32]{Sun2010v11}. We specify the
version because later versions of that preprint reorganized the
conjectures. The all-index statement includes every composite
integer and every prime power; the second statement includes all
primes, without an exceptional prime.

The principal arithmetic input for the all-index proof is He's
published cubic-binomial divisibility \cite{He2017}, stated explicitly
in Theorem~\ref{s57:he-input}. All formal-coordinate, valuation,
interpolation, and propagation steps needed to pass from that input
to \eqref{eq:main-allN} are proved below. No complex multiplication,
ordinary-prime hypothesis, numerical assumption, or conjectural
supercongruence is used in that argument.

The prime-square assertion has a different status in the proof.
We include a finite, self-contained proof of its arithmetic reduction
and its cubic-binomial congruence. However, the cubic congruence was
already known in stronger forms \cite{SunEuler2011}, and the
weighted reduction can also be obtained from the published
Ap\'ery-polynomial results of Sun and Guo--Zeng
\cite{SunApery2012,GuoZeng2012}. Thus the prime-square calculation
is not presented as a historically new arithmetic principle.

\section{A constant-term identity for the original sum}
\label{sec:constant-term}

All formal series in this section have integral coefficients. Put
\[
 \Lambda=\frac{(1+u)(1+v)}{uv},\qquad r=-\Lambda,
 \qquad \rho=\frac{uv}{1+u+v+2uv},
 \qquad G(x)=\sum_{l\ge0}(-16)^l b_l^2x^{2l}.
\]
Here \(\CT\) denotes the coefficient of \(u^0v^0\). Products of
a Laurent polynomial with the displayed formal series have
well-defined integral constant terms. Direct extraction gives
\[
                  w_n=\CT\bigl(\Lambda^nG(uv)\bigr).
\]
Since \((r-1)^{-1}=-\rho\), the finite geometric-sum identity and
its formal derivative give
\begin{align}
 S(N)&=\CT\bigl(r^N\Psi\bigr)
          -42N\CT\bigl(r^N\rho G(uv)\bigr),
          \label{eq:signed-CT}\\
 \Psi&=(5\rho-42\rho^2)G(uv).
          \label{eq:signed-Psi}
\end{align}
For clarity, for arbitrary \(A,B\) one has
\[
 \sum_{n=0}^{N-1}(An+B)r^n
 =\frac{ANr^N}{r-1}
  +(r^N-1)\left(\frac B{r-1}-\frac{Ar}{(r-1)^2}\right).
\]
Substitute \(A=42,B=37\) and
\(r/(r-1)^2=-\rho+\rho^2\). The terms without \(r^N\) have
zero constant term after multiplication by \(G(uv)\), proving
\eqref{eq:signed-CT}.

Let \(\theta_u=u\partial_u\) and \(\theta_v=v\partial_v\).
If
\begin{equation}\label{eq:divergence-goal}
             \Psi=\theta_u P+\theta_v Q,
                 \qquad P,Q\in\Z[[u,v]],
\end{equation}
then formal integration by parts shows that
\[
 \CT(r^N\Psi)
 =-N\CT\left(r^{N-1}
       \bigl((\theta_u r)P+(\theta_v r)Q\bigr)\right)\in N\Z.
\]
Together with \eqref{eq:signed-CT}, this proves
\eqref{eq:main-allN}. The next two sections establish
\eqref{eq:divergence-goal} with integral coefficients.

% Integral formal-series lemma, extending the alpha=0 Catalan criterion.
% No analytic convergence or prime restriction is involved.

\section{A diagonal criterion with a bilinear denominator}

Put $A=\mathbb Z[[u,v]]$, $\theta_u=u\partial_u$,
$\theta_v=v\partial_v$, and
$\mathcal D_A=\theta_u A+\theta_v A$.

\begin{theorem}[Integral diagonal criterion]\label{thm:diag-criterion}
Let $\kappa,a,b\in\mathbb Z$, let $R(x)\in\mathbb Z[[x]]$, and put
\[
 \rho_\kappa=\frac{uv}{1+u+v+\kappa uv},\qquad
 F(u,v)=(a\rho_\kappa+b\rho_\kappa^2)R(uv).
\]
Then
\[
 F\in\mathcal D_A
 \quad\Longleftrightarrow\quad
 n\mid[u^nv^n]F\quad\text{for every }n\ge1.
\]
\end{theorem}

\begin{proof}
For any $F\in A$, coefficientwise Bezout gives
\[
 F\in\mathcal D_A
 \quad\Longleftrightarrow\quad
 [u^0v^0]F=0\quad\text{and}\quad
 \gcd(r,s)\mid[u^rv^s]F\quad((r,s)\ne(0,0)).
\]
The forward implication of the theorem follows immediately.

For the converse, let $C_\kappa(x)\in1+x\mathbb Z[[x]]$ be
the unique solution of
\[
 C_\kappa=1+xC_\kappa^2-\kappa xC_\kappa.
\]
This recursion proves existence, uniqueness and integrality
coefficient by coefficient. Define
\[
 U=uC_\kappa(uv),\qquad V=vC_\kappa(uv),\qquad y=UV.
\]
This is an integral formal automorphism. More explicitly, let
$B(y)\in1+y\mathbb Z[[y]]$ be the unique solution of
\[
 B^2-(1+y)B+\kappa y=0,
 \qquad\text{equivalently}\qquad
 B=1+y-\frac{\kappa y}{B}.
\]
Again the latter equation proves integrality and uniqueness by
recursion. The inverse coordinate change is
\[
 u=\frac U{B(y)},\qquad v=\frac V{B(y)},\qquad
 x=uv=\frac y{B(y)^2},\qquad C_\kappa(x)=B(y).
\]
Write $e(y)=yB'(y)/B(y)$. The inverse logarithmic differential
matrix and its determinant are
\[
 \begin{pmatrix}d\log u\\d\log v\end{pmatrix}
 =\begin{pmatrix}1-e&-e\\-e&1-e\end{pmatrix}
 \begin{pmatrix}d\log U\\d\log V\end{pmatrix},
 \qquad J(y)=1-2e(y)\in1+y\mathbb Z[[y]].
\]
Both this matrix and its inverse have integral power-series
entries. Integral exact logarithmic two-forms are therefore
preserved in both directions. In other words, $F\in\mathcal D_A$
if and only if
\[
 \widetilde F(U,V)=J(y)F\left(\frac U{B(y)},\frac V{B(y)}\right)
\]
is an integral logarithmic divergence in $U,V$.

The defining equation for $B$ gives the factorization
\[
 1+u+v+\kappa uv=\frac{(1+U)(1+V)}{B(y)},
 \qquad
 \rho_\kappa=\frac{y}{B(y)(1+U)(1+V)}.
\]
Consequently
\[
 \widetilde F
 =\frac{A_1(y)}{(1+U)(1+V)}
  +\frac{A_2(y)}{(1+U)^2(1+V)^2},
\]
where
\[
 A_1(y)=\frac{aJ(y)y}{B(y)}R\left(\frac y{B(y)^2}\right),
 \qquad
 A_2(y)=\frac{bJ(y)y^2}{B(y)^2}R\left(\frac y{B(y)^2}\right).
\]
These are integral and start in degrees one and two respectively.
Write $A_j(y)=\sum_k A_{j,k}y^k$. For $r,s\ge1$, put
$q=\min(r,s)$ and $g=\gcd(r,s)$. Then
\[
 [U^rV^s]\widetilde F
 =(-1)^{r+s}\sum_{k=0}^{q}
 \bigl(A_{1,k}+A_{2,k}(r-k+1)(s-k+1)\bigr).
\]
Since $r,s,q$ are multiples of $g$, this yields
\[
 [U^rV^s]\widetilde F
 \equiv(-1)^{r+s}[U^qV^q]\widetilde F\pmod g.
\]

It remains to check the diagonal divisibility after substitution.
Write $\operatorname{Diag}F(x)=\sum_n[u^nv^n]F\,x^n$.
The radial substitution preserves the difference of the two
exponents, so
\[
 \operatorname{Diag}\widetilde F(y)
 =J(y)\operatorname{Diag}F\left(\frac y{B(y)^2}\right).
\]
By the hypothesis there is $H_0(x)\in x\mathbb Z[[x]]$ such that
$\operatorname{Diag}F(x)=xH_0'(x)$. For $x(y)=y/B(y)^2$,
\[
 \frac{yx'(y)}{x(y)}=1-2\frac{yB'(y)}{B(y)}=J(y),
\]
and hence
\[
 \operatorname{Diag}\widetilde F(y)
 =y\frac d{dy}H_0(x(y)).
\]
Its coefficient of $y^q$ is divisible by $q$. The preceding
coefficient congruence now gives all gcd conditions. Constant
and axis coefficients vanish, so the Bezout criterion proves
$\widetilde F$ is an integral logarithmic divergence. Pulling
back proves the same for $F$.
\end{proof}

\begin{lemma}[The diagonal generating function]
With the preceding notation, put
\[
 Q_\kappa(x)=(1+\kappa x)^2-4x.
\]
Then
\[
 \operatorname{Diag}F(x)
 =\frac{xR(x)\bigl(aQ_\kappa(x)+bx(1+\kappa x)\bigr)}
        {Q_\kappa(x)^{3/2}}.
\]
\end{lemma}

\begin{proof}
Expanding geometrically in $u+v$ gives
\[
 \operatorname{Diag}\frac1{1+u+v+\kappa uv}
 =\sum_{i\ge0}\binom{2i}i\frac{x^i}{(1+\kappa x)^{2i+1}}
 =Q_\kappa(x)^{-1/2}.
\]
The analogous expansion for the squared denominator gives
$(1+\kappa x)Q_\kappa(x)^{-3/2}$. Multiplication by $uv$,
$(uv)^2$ and $R(uv)$ yields the formula.
\end{proof}

\begin{corollary}[The signed kernel for Conjecture 5.7]\label{cor:signed-kernel}
Let
\[
 G(x)=\sum_{\ell\ge0}(-16)^\ell\binom{2\ell}\ell^2x^{2\ell},
 \quad \rho=\frac{uv}{1+u+v+2uv},
 \quad\Psi=(5\rho-42\rho^2)G(uv).
\]
Then
\[
 \Psi\in\theta_u\mathbb Z[[u,v]]+\theta_v\mathbb Z[[u,v]]
 \quad\Longleftrightarrow\quad
 N\mid[x^{N-1}]\frac{G(x)(5-42x-64x^2)}{(1+4x^2)^{3/2}}
 \quad(N\ge1).
\]
\end{corollary}

% Arithmetic fragment for Sun Conjecture 5.7(ii), all-N assertion.
% No preamble. Required packages: amsmath, amssymb, amsthm.
% Required theorem environments: theorem, lemma, proposition, corollary.
% Required bibliography key: He2017.

\section{The integral diagonal}
\label{s57:arithmetic}

Write \(b_j=\binom{2j}{j}\). The diagonal sequence needed for the
signed constant-term reduction is
\begin{align}
 Y(2h+1)
   &=\sum_{i+l=h}(42i+5)(-1)^i(-16)^l b_i b_l^2,
       &&h\ge0,\label{s57:odd-diagonal}\\
 Y(2M)
   &=-42\sum_{i+l=M-1}(2i+1)(-1)^i(-16)^l b_i b_l^2,
       &&M\ge1.\label{s57:even-diagonal}
\end{align}
All indices in these sums are nonnegative integers.
We prove the complete divisibility statement
\[
                         N\mid Y(N)\qquad(N\ge1).
\]

\subsection{An established Ramanujan partial-sum divisibility}

The one arithmetic input is the following published theorem. It is
not claimed as a new result here.

\begin{theorem}[Bing He \cite{He2017}]
\label{s57:he-input}
For every integer \(n>1\),
\[
 2n\binom{2n}{n}\ \Bigm|\
 \sum_{k=0}^{n-1}(42k+5)\binom{2k}{k}^{\!3}4096^{\,n-1-k}.
\]
\end{theorem}

Define
\[
 A_j=\frac{(42j+5)b_j^3}{4096^j},\qquad
 R_j=\sum_{k=0}^j A_k.
\]
For every odd prime \(p\) and every \(j\ge0\),
\begin{equation}
 v_p(R_j)\ge v_p\bigl((2j+1)(2j+2)\bigr).
 \label{s57:partial-bound}
\end{equation}
Indeed, for \(j\ge1\), apply Theorem~\ref{s57:he-input} with
\(n=j+1\) and use the exact equality
\[
 2(j+1)b_{j+1}=4(2j+1)b_j.
\]
The denominator \(4096^j\) is a \(p\)-adic unit, and the Catalan
integrality \(b_j/(j+1)\in\mathbb Z\) gives the bound.
For \(j=0\), it follows from \(R_0=5\) and the fact that \(p\) is odd.

\subsection{Every odd prime-power base}

\begin{theorem}
\label{s57:prime-power-bases}
For every odd prime \(p\) and every integer \(e\ge1\),
\[
                            p^e\mid Y(p^e).
\]
\end{theorem}
\begin{proof}
Put \(N=p^e\), \(h=(N-1)/2\), and
\[
 r_j=\prod_{a=1}^j
       \left(\frac{1-N/(2a-1)}{1-N/(2a)}\right)^2
       \quad(0\le j\le h),\qquad r_0=1.
\]
The factors are \(p\)-adic units: all positive denominators
\(2a-1,2a\) are less than \(N\), so their valuations are smaller
than \(e\). Comparing consecutive ratios, starting at \(j=0\), gives
the exact identity
\[
 \left(\frac{b_{h-j}}{b_h}\right)^2
                  =\frac{b_j^2}{256^j}\,r_j.
\]
Substitution \(l=h-j\), \(i=j\) into
\eqref{s57:odd-diagonal} therefore gives
\begin{equation}
 Y(N)=(-16)^h b_h^2\sum_{j=0}^h A_jr_j.
 \label{s57:abel-bridge}
\end{equation}

For \(0\le j<h\), set \(a=2j+1\), \(b=2j+2\). Direct subtraction
of the consecutive products gives
\[
 r_{j+1}-r_j
  =-N r_j\frac{2ab-N(a+b)}{a^2(b-N)^2}.
\]
The integers \(a,b\) are consecutive and smaller than \(N\).
If \(t=v_p(ab)\), then \(t<e\), \(v_p(b-N)=v_p(b)\), and
\(v_p(2ab-N(a+b))=t\). The last equality uses that \(p\) is odd:
the first term has valuation \(t\), and the second has valuation
at least \(e\). It follows that
\[
 v_p(r_{j+1}-r_j)
       =e-v_p\bigl((2j+1)(2j+2)\bigr).
\]
Abel summation now yields
\[
 \sum_{j=0}^h A_jr_j
   =R_hr_h-\sum_{j=0}^{h-1}R_j(r_{j+1}-r_j)\in p^e\mathbb Z_p.
\]
Here each term of the last sum has valuation at least \(e\) by
\eqref{s57:partial-bound}, while \(R_h\in p^e\mathbb Z_p\) because
\(2h+1=N\). The integer prefactor in \eqref{s57:abel-bridge} proves
the theorem. This argument includes \(p=3\) and \(p=7\).
\end{proof}

\subsection{Integral central-binomial digit kernels}

We give the integral interpolation needed to propagate these bases;
the proof does not discard carries. Fix an arbitrary odd prime \(p\).
The ring \(\mathbb Z_p\langle X_1,\ldots,X_d\rangle\) consists of
power series whose coefficients lie in \(\mathbb Z_p\) and tend to
zero as total degree tends to infinity. It is complete for the Gauss
norm. Integral affine substitutions preserve the ring, and evaluation
at points of \(\mathbb Z_p^d\) is well defined.

\begin{lemma}[Normalized unit factorials]
\label{s57:unit-factorials}
Put
\[
 E_p(M)=\prod_{\substack{1\le j\le M\\p\nmid j}}j,
 \qquad W=(p-1)!.
\]
For each \(0\le r<p\) there is a unit
\(\mathcal H_r(X)\in1+p\mathbb Z_p\langle X\rangle\), with
\(\mathcal H_r(0)=1\), such that
\[
             E_p(pq+r)=W^q r!\,\mathcal H_r(q)
                  \qquad(q\ge0).
\]
\end{lemma}
\begin{proof}
Set
\[
 B(X)=\prod_{i=1}^{p-1}(1+pX/i),\qquad
 \log B(X)=\sum_{j\ge1}\ell_jX^j,
 \quad
 \ell_j=\frac{(-1)^{j-1}p^j}{j}\sum_{i=1}^{p-1}i^{-j}.
\]
The polynomial continuation of \(\sum_{q=0}^{n-1}q^j\) is
\[
 S_j(X)=\sum_{a=0}^j
 \left\{\begin{matrix}j\\a\end{matrix}\right\}
 \frac{(X)_{\underline{a+1}}}{a+1},
\]
where the braces denote Stirling numbers of the second kind and the
underlined subscript denotes a falling factorial. In particular,
\(S_j\in p^{-\lfloor\log_p(j+1)\rfloor}\mathbb Z_p[X]\).
For every odd \(p\) and \(j\ge1\),
\[
 j-v_p(j)-\lfloor\log_p(j+1)\rfloor\ge1,
\]
and these bounds tend to infinity with \(j\). The inequality is
immediate for \(j=1,2,3\). For \(j\ge4\), use
\(v_p(j)\le\lfloor\log_3(j+1)\rfloor\) and
\(2\lfloor\log_3(j+1)\rfloor\le j-1\); the latter follows directly
for \(4\le j\le7\), and then from
\(2t\le3^t-2\) when \(t=\lfloor\log_3(j+1)\rfloor\ge2\).
Consequently,
\[
 A(X)=\sum_{j\ge1}\ell_jS_j(X)
                   \in p\mathbb Z_p\langle X\rangle.
\]
Since \(p\) is odd, \(\exp A(X)\) converges in
\(1+p\mathbb Z_p\langle X\rangle\): its \(k\)-th term has valuation
at least \(k-v_p(k!)\), which tends to infinity. Define
\[
 \mathcal H_r(X)=\exp A(X)\prod_{i=1}^r(1+pX/i).
\]
All \(S_j(0)\) vanish, giving \(\mathcal H_r(0)=1\). For every
nonnegative integer \(q\), convergence gives
\[
 A(q)=\sum_{a=0}^{q-1}\log B(a),\qquad
 \exp A(q)=\prod_{a=0}^{q-1}B(a).
\]
Splitting \(E_p(pq+r)\) into its complete blocks and its last \(r\)
factors proves the assertion. Only the normalized factor is
interpolated; no interpolation of \(W^q\) is needed.
\end{proof}

\begin{lemma}[Central-binomial kernel]
\label{s57:central-kernel}
For \(0\le r<p\), put
\[
 \delta_r=\left\lfloor\frac{2r}{p}\right\rfloor,
 \qquad t_r=2r-p\delta_r,
\]
and define
\[
 \mathcal B_r(X)=
 p^{\delta_r}(2X+1)^{\delta_r}W^{\delta_r}
 \frac{t_r!}{r!^2}
 \frac{\mathcal H_{t_r}(2X+\delta_r)}{\mathcal H_r(X)^2}.
\]
Then \(\mathcal B_r\in p^{\delta_r}\mathbb Z_p\langle X\rangle\),
and for all \(q\ge0\),
\[
 b_{pq+r}=b_q\mathcal B_r(q),\qquad
 \mathcal B_r(0)=b_r,\qquad
 \mathcal B_r(X)\equiv b_r\pmod p.
\]
\end{lemma}
\begin{proof}
All factorial denominators displayed in \(\mathcal B_r\) are units.
Use the identity
\[
 M!=p^{\lfloor M/p\rfloor}\lfloor M/p\rfloor!E_p(M)
\]
in \(b_{pq+r}\), followed by Lemma~\ref{s57:unit-factorials}.
The factorial part remaining after removing the unit factorials is
\[
 p^{\delta_r}\frac{(2q+\delta_r)!}{q!^2}
       =p^{\delta_r}b_q(2q+1)^{\delta_r},
\]
and the powers of \(W\) reduce to \(W^{\delta_r}\). This proves the
exact formula. Evaluation at \(q=0\) gives the constant value.
If \(\delta_r=1\), both sides of the asserted congruence vanish
modulo \(p\); if \(\delta_r=0\), all \(\mathcal H\)'s are one
modulo \(p\), and \(t_r=2r\), proving the congruence.
\end{proof}

Put
\[
 w(i,l)=(-1)^i(-16)^l b_i b_l^2,
 \qquad \mathcal A_p=\mathbb Z_p\langle X,Y\rangle.
\]
The restricted series
\[
 \mathcal C(Y)=\exp\bigl(Y\log((-16)^{p-1})\bigr)
                    \in1+p\mathbb Z_p\langle Y\rangle
\]
satisfies \(\mathcal C(j)=(-16)^{(p-1)j}\) for integers \(j\ge0\).
Define
\begin{equation}
 K_{r,s}(X,Y)=(-1)^r(-16)^s\mathcal C(Y)
                  \mathcal B_r(X)\mathcal B_s(Y)^2.
 \label{s57:two-variable-kernel}
\end{equation}
Because \(p\) is odd, \((-1)^{pq+r}=(-1)^{q+r}\), so
Lemma~\ref{s57:central-kernel} gives the exact factorization
\begin{equation}
 w(pq+r,pj+s)=w(q,j)K_{r,s}(q,j)\qquad(q,j\ge0).
 \label{s57:exact-kernel}
\end{equation}
It also gives
\begin{equation}
 \begin{gathered}
 K_{r,s}\in p^{\delta_r+2\delta_s}\mathcal A_p,
 \qquad K_{r,s}(0,0)=w(r,s),\\
 K_{r,s}(X,Y)\equiv w(r,s)\pmod p.
 \end{gathered}
 \label{s57:kernel-properties}
\end{equation}

\subsection{A two-state operator on odd indices}

For \(F=(F_0,F_1)\in\mathcal M_p:=\mathcal A_p^2\) and every
positive odd integer \(N\), define
\[
 \mathcal J_N(F)=
 \sum_{i+l=(N-1)/2}w(i,l)F_0(i,l)
 +\sum_{i+l=(N-1)/2-1}w(i,l)F_1(i,l).
\]
A sum with a negative right-hand side is empty; in particular,
\(\mathcal J_1(F)=F_0(0,0)\). Let \(h_p=(p-1)/2\).

\begin{proposition}[Exact digit decomposition]
\label{s57:digit-operator}
For \(d=0,1\), set
\[
 (\mathcal L_pF)_d(X,Y)=
 \sum_{a=0}^1
 \sum_{\substack{0\le r,s<p\\r+s=h_p-a+pd}}
 K_{r,s}(X,Y)F_a(pX+r,pY+s).
\]
Then \(\mathcal L_p\) maps \(\mathcal M_p\) to itself, and for
every positive odd integer \(N\),
\[
                   \mathcal J_{pN}(F)=\mathcal J_N(\mathcal L_pF).
\]
Moreover, there is \(R_F\in\mathcal M_p\) with
\[
 \mathcal L_pF=(\mathcal J_p(F),0)+pR_F,
 \qquad (R_F)_0(0,0)=0.
\]
\end{proposition}
\begin{proof}
In the input state \(a\in\{0,1\}\), write \(i=pq+r\),
\(l=pj+s\), \(0\le r,s<p\). The constraint becomes
\[
 p(q+j)+r+s
     =\frac{pN-1}{2}-a
     =p\frac{N-1}{2}+h_p-a.
\]
Thus
\[
 r+s=h_p-a+pd,\qquad q+j=\frac{N-1}{2}-d.
\]
Only \(d=0,1\) is possible: \(h_p-a\ge0\), and
\(0\le r+s\le2p-2\). This includes \(p=3\), when the smallest
value \(h_p-a\) is zero. Equation~\eqref{s57:exact-kernel} proves
the exact block identity, with all empty ranges respected. Integral
affine substitution and the finite sum preserve \(\mathcal M_p\).

In output state \(d=1\),
\[
 r+s=p+h_p-a\ge p+h_p-1\ge p>p-1.
\]
If neither central binomial coefficient carried, then
\(r,s\le h_p\), contradicting this inequality. Hence
\(\delta_r+2\delta_s\ge1\), and
\((\mathcal L_pF)_1\in p\mathcal A_p\).

For \(d=0\), use \eqref{s57:kernel-properties} and
\[
 F_a(pX+r,pY+s)\equiv F_a(r,s)\pmod p.
\]
The latter holds first for polynomials and then for restricted power
series by Gauss-norm continuity. The resulting constant is exactly
\(\mathcal J_p(F)\). Moreover, evaluating the first component at
\((X,Y)=(0,0)\) gives this constant exactly, because
\(K_{r,s}(0,0)=w(r,s)\). Subtracting the constant vector and
dividing by \(p\) proves the last assertion, including its zero
constant term.
\end{proof}

\begin{theorem}[Propagation from prime-power bases]
\label{s57:propagation}
Let \(p\) be an odd prime, \(F\in\mathcal M_p\), and \(e\ge1\).
Suppose that
\[
            \mathcal J_{p^j}(F)\in p^j\mathbb Z_p
                    \qquad(1\le j<e).
\]
Then, for every positive odd integer \(m\),
\[
 \mathcal J_{p^em}(F)
    \equiv \mathcal J_{p^e}(F)\mathcal J_m(1,0)\pmod{p^e}.
\]
If the base divisibility also holds at \(j=e\), then
\(\mathcal J_{p^em}(F)\in p^e\mathbb Z_p\) for every such \(m\).
No coprimality condition on \(m\) is required.
\end{theorem}
\begin{proof}
Iterating the exact block identity gives
\[
                  (\mathcal L_p^jF)_0(0,0)=\mathcal J_{p^j}(F).
\]
Suppose the first \(j-1\) base divisibilities have given
\(\mathcal L_p^{j-1}F=p^{j-1}G\), \(G\in\mathcal M_p\); for
\(j=1\), take \(G=F\). Proposition~\ref{s57:digit-operator} yields
\[
 \mathcal L_p^jF
   =p^{j-1}(\mathcal J_p(G),0)+p^jR_G
   =(\mathcal J_{p^j}(F),0)+p^jR_G.
\]
The second equality follows either from the exact block identity or
by evaluation at \((0,0)\), using \((R_G)_0(0,0)=0\).
The next base divisibility therefore puts the whole iterate in
\(p^j\mathcal M_p\). Induction through \(e-1\), and then one
more application of the formula, gives
\[
       \mathcal L_p^eF=(\mathcal J_{p^e}(F),0)+p^eR
                    \qquad(R\in\mathcal M_p).
\]
Apply \(\mathcal J_m\), which is integral on \(\mathcal M_p\).
The hypothesis is empty for \(e=1\), and the same argument applies.
The case \(m=1\) is included throughout.
\end{proof}

\begin{corollary}
\label{s57:odd-diagonal-integrality}
For every positive odd integer \(N\), one has \(N\mid Y(N)\).
\end{corollary}
\begin{proof}
Take \(F=(42X+5,0)\), so \(\mathcal J_N(F)=Y(N)\) for odd
\(N\). Theorem~\ref{s57:prime-power-bases} supplies every base
required by Theorem~\ref{s57:propagation}. For any prime power
\(p^e\mid N\), the cofactor \(m=N/p^e\) is odd; hence
\(p^e\mid Y(N)\). This proves the assertion for all odd \(N>1\).
For \(N=1\), it is immediate from \(Y(1)=5\). Exponent zero
requires no divisibility condition.
\end{proof}

\subsection{Even indices: an integral individual summand}

\begin{lemma}
\label{s57:termwise-even}
For all nonnegative integers \(i,l\),
\[
             i+l+1\ \Bigm|\ (2i+1)b_i b_l.
\]
\end{lemma}
\begin{proof}
Set \(M=i+l+1\), and fix \(p^e\mid M\) with \(e\ge1\).
For \(1\le j\le e\), put \(q=p^j\), and let \(r,s\) be the
residues of \(i,l\) modulo \(q\), respectively. Then
\(r+s=q-1\): their sum is congruent to \(-1\) modulo \(q\),
and lies between zero and \(2q-2\).
The contribution of this level to \(v_p(b_i)+v_p(b_l)\) is
\[
                  \left\lfloor\frac{2r}{q}\right\rfloor
                  +\left\lfloor\frac{2s}{q}\right\rfloor.
\]
If \(q\) is even, this equals one. If \(q\) is odd, it also equals
one unless \(r=s=(q-1)/2\); in that exceptional case it is zero,
but \(q\mid2i+1\). Legendre's factorial valuation formula and
\(v_p(2i+1)=\sum_{j\ge1}\mathbf 1_{p^j\mid2i+1}\) show that
the total valuation of \((2i+1)b_i b_l\) is at least \(e\).
This holds for every prime power dividing \(M\), proving the lemma.
\end{proof}

\begin{theorem}[Complete diagonal integrality]
\label{s57:complete-diagonal}
For every integer \(N\ge1\),
\[
                              N\mid Y(N).
\]
\end{theorem}
\begin{proof}
Odd indices are covered by
Corollary~\ref{s57:odd-diagonal-integrality}. For \(N=2M\),
each summand in \eqref{s57:even-diagonal} is divisible by \(2M\):
Lemma~\ref{s57:termwise-even} supplies the factor \(M\), the
remaining factor \(b_l\) and all powers are integers, and the
coefficient \(42\) supplies the factor two. This includes \(M=1\).
\end{proof}


\section{Completion of the all-index assertion}
\label{sec:finish-allN}

The odd and even formulas \eqref{s57:odd-diagonal} and
\eqref{s57:even-diagonal} are the coefficients of the diagonal
generating function in Corollary~\ref{cor:signed-kernel}.
Indeed,
\[
 (1+4z)^{-3/2}=\sum_{i\ge0}(-1)^i(2i+1)b_i z^i,
\]
and
\[
 [z^i]\frac{5-64z}{(1+4z)^{3/2}}
                =(-1)^i(42i+5)b_i.
\]
For \(i\ge1\), the latter identity uses
\(b_{i-1}/b_i=i/(2(2i-1))\); for \(i=0\), both sides are five.
The odd term \(-42x\) in the numerator gives exactly
\eqref{s57:even-diagonal}.

Theorem~\ref{s57:complete-diagonal} proves every diagonal
divisibility required by Corollary~\ref{cor:signed-kernel}.
The integral diagonal criterion therefore supplies \(P,Q\)
in \eqref{eq:divergence-goal}. Applying the integration-by-parts
identity of Section~\ref{sec:constant-term} proves
\eqref{eq:main-allN} for every \(N\ge1\).
This completes the first assertion of Theorem~\ref{thm:main}.

% Self-contained proof fragment for the second assertion of Sun 5.7(ii).
% Requires amsmath, amssymb, amsthm and theorem/lemma environments.
% All mathematical ingredients used below are proved in this fragment.
% Bibliographic keys to supply in the containing paper:
% SunApery2012 = Z.-W. Sun, On sums of Apery polynomials and related
% congruences, J. Number Theory 132 (2012), 2673--2699; arXiv:1101.1946v4.
% GuoZeng2012 = V. J. W. Guo and J. Zeng, Proof of some conjectures of
% Z.-W. Sun on congruences for Apery polynomials, J. Number Theory 132
% (2012), 1731--1740; arXiv:1101.0983v2.
% Guillera2008 = J. Guillera, arXiv:1104.0396, Identity 2 and its second
% WZ pair. The exact certificate below is independently verified.
% SunEuler2011 = Z.-W. Sun, Super congruences and Euler numbers,
% Science China Mathematics 54 (2011), no. 12, 2509--2535;
% arXiv:1001.4453v20, Theorem 1.3, equation (1.15).

\section{Sun's signed weighted congruence modulo a prime square}

Put $b_k=\binom{2k}k$ and
\[
 W_n(1,t)=\sum_{k=0}^{\lfloor n/2\rfloor}
 t^k b_k^2\binom n{2k}^{\!2}.
\]
We prove the unchanged prime-square assertion of
Conjecture~5.7(ii): for every prime $p$,
\begin{equation}\label{eq:sun57-target}
 \sum_{n=0}^{p-1}(42n+37)(-1)^nW_n(1,-16)
 \equiv p\left(21\left(\frac p7\right)+16\right)\pmod{p^2}.
\end{equation}
Here $(p/7)$ is the Legendre symbol, with value zero when $p=7$.

Sun's unweighted transformation and Guo--Zeng's weighted
Ap\'ery-polynomial congruence \cite{SunApery2012,GuoZeng2012}
also give the signed transfer by an elementary index shift and
coefficientwise differentiation. We give a direct alternative
derivation through an alternating polynomial functional. The cubic
auxiliary congruence used here follows from the stronger result
of Sun \cite[Theorem 1.3, (1.15)]{SunEuler2011};
we include a finite algebraic proof to keep the argument independent
of any supercongruence input. None of these observations establishes
historical priority for the original $W$-congruence.

\subsection{An alternating polynomial functional}

For $f\in\mathbb Q[x]$, let $Q_f\in\mathbb Q[x]$ be the unique
polynomial satisfying
\[
 Q_f(x+1)+Q_f(x)=2f(x).
\]
The operator $Q\mapsto Q(x+1)+Q(x)$ acts triangularly on monomials,
with diagonal entries $2$, proving existence and uniqueness. If
$p$ is odd and $f\in\mathbb Z_{(p)}[x]$, the same triangular
argument shows $Q_f\in\mathbb Z_{(p)}[x]$. For every odd positive
integer $N$, telescoping gives the exact identity
\begin{equation}\label{eq:alternating-functional}
 \sum_{n=0}^{N-1}(-1)^n f(n)=\frac{Q_f(0)+Q_f(N)}2.
\end{equation}
In particular, for a $p$-integral polynomial $f$ and an odd prime $p$,
the right side is
$Q_f(0)+(p/2)Q_f'(0)$ modulo $p^2$.

\begin{lemma}\label{lem:alternating-binomial-functional}
Let $f_m(x)=\binom xm^2$ and $Q_m=Q_{f_m}$. For every $k\ge0$,
\[
 Q_{2k}(0)=\frac{(-1)^k b_k}{16^k},\qquad
 Q'_{2k}(0)=0,\qquad
 Q'_{2k+1}(0)=\frac{(-1)^{k+1}}{(2k+1)b_k}.
\]
\end{lemma}

\begin{proof}
For a polynomial $f$, the series
$F_f(z)=\sum_{n\ge0}f(n)z^n$ is a rational function whose only
possible finite pole is at $z=1$. Summing the defining relation for
$Q_f$ gives
\[
 (1+z)F_{Q_f}(z)-Q_f(0)=2zF_f(z).
\]
Evaluation at $z=-1$ therefore yields $Q_f(0)=2F_f(-1)$.
The elementary binomial generating identity
\[
 \sum_{n\ge0}\binom nm^2z^n
 =\frac{z^m}{(1-z)^{2m+1}}
    \sum_{j=0}^m\binom mj^2z^j
\]
follows from binomial product linearization and the binomial
theorem. At $m=2k$, use
\[
 \sum_{j=0}^{2k}(-1)^j\binom{2k}j^2=(-1)^k\binom{2k}k,
\]
obtained by taking the coefficient of $X^{2k}$ in
$(1-X)^{2k}(1+X)^{2k}$. The first asserted formula follows.

The identity $f_m(m-1-x)=f_m(x)$ and uniqueness imply
$Q_m(m-x)=Q_m(x)$. The double zeros of $f_m$ at
$0,1,\ldots,m-1$ give, for $0\le j\le m$,
\[
 Q_m(j)=(-1)^jQ_m(0),\qquad
 Q'_m(j)=(-1)^jQ'_m(0).
\]
For even $m$, differentiation of the reflection identity gives
$Q'_m(m)=-Q'_m(0)$, while the recurrence gives equality.
Thus $Q'_m(0)=0$. The case $m=0$ is immediate.

For odd $m=2k+1$, reflection of the values gives $Q_m(0)=0$,
and consequently $Q_m$ vanishes at all $j=0,\ldots,m$. Put
$P(x)=\prod_{j=0}^m(x-j)$. The proper rational function
$Q_m(x)/P(x)^2$ has only simple poles, with residues
\[
 \frac{Q'_m(j)}{P'(j)^2}
 =\frac{(-1)^jQ'_m(0)}{j!^2(m-j)!^2}.
\]
Since $Q_m$ has degree $2m$ and leading coefficient $1/m!^2$,
the coefficient of $x^{-2}$ at infinity is $1/m!^2$. Comparing
with the partial-fraction expansion yields
\[
 1=Q'_m(0)\sum_{j=0}^m(-1)^j j\binom mj^2.
\]
For $m=2k+1$, the sum on the right is
$(-1)^{k+1}(2k+1)b_k$. To see this, use
$j\binom mj=m\binom{m-1}{j-1}$ and take the coefficient of
$X^{m-1}$ in
$(1-X)^{m-1}(1+X)^m=(1-X^2)^{2k}(1+X)$.
This proves the remaining derivative formula.
\end{proof}

\subsection{The two signed inner sums}

\begin{lemma}\label{lem:signed-inner-sums}
Let $p$ be odd and $0\le2k<p$. Then, in $\mathbb Z_{(p)}$,
\begin{align*}
 \sum_{n=0}^{p-1}(-1)^n\binom n{2k}^{\!2}
 &\equiv\frac{(-1)^k b_k}{16^k}\pmod{p^2},\\
 \sum_{n=0}^{p-1}n(-1)^n\binom n{2k}^{\!2}
 &\equiv\left(k-\frac12\right)\frac{(-1)^k b_k}{16^k}
       +\frac p2\frac{(-1)^k}{b_k}\pmod{p^2}.
\end{align*}
\end{lemma}

\begin{proof}
The polynomial $f_{2k}$ has $p$-integral coefficients because
$(2k)!$ is a unit. Its value in
\eqref{eq:alternating-functional}, together with
Lemma~\ref{lem:alternating-binomial-functional}, proves the first
congruence.

For the second, put $M_m=Q_{x f_m}$. The polynomial identity
\[
 (2x-m+1)f_m(x)
 =(m+1)\bigl(f_{m+1}(x+1)-f_{m+1}(x)\bigr)
\]
and the defining equation for $Q$ give
\[
 2M_m-(m-1)Q_m=2(m+1)(f_{m+1}-Q_{m+1}).
\]
At $m=2k$, one has $Q_{m+1}(0)=0$ and
$f_{m+1}(0)=f'_{m+1}(0)=0$. Hence
\[
 M_{2k}(0)=\left(k-\frac12\right)Q_{2k}(0),
 \qquad M'_{2k}(0)=\frac{(-1)^k}{b_k}.
\]
Apply \eqref{eq:alternating-functional} to $xf_{2k}$.

This proof includes the endpoint $2k=p-1$. Although the rational
derivation passes through $f_{2k+1}=f_p$ there, coefficientwise
$p$-integrality is needed only for $f_{2k}$ and $xf_{2k}$, whose
denominators are units. Also $b_k$ is a unit whenever $2k<p$.
\end{proof}

\begin{theorem}[General signed transfer]\label{thm:signed-transfer}
Let $p$ be an odd prime, $h=(p-1)/2$, and $A,B,t\in\mathbb Z$.
Then
\begin{align}
 \sum_{n=0}^{p-1}(An+B)(-1)^nW_n(1,t)
 &\equiv\sum_{k=0}^h
 \left(Ak+B-\frac A2\right)
 \left(-\frac t{16}\right)^k b_k^3\notag\\
 &\hspace{1em}+\frac A2p\left(\frac{1+4t}{p}\right)
 \pmod{p^2}.
 \label{eq:general-signed-transfer}
\end{align}
\end{theorem}

\begin{proof}
Interchange the finite sums and apply
Lemma~\ref{lem:signed-inner-sums}. The remaining term is
\[
 \frac A2p\sum_{k=0}^h(-t)^k b_k.
\]
Since $b_k\equiv(-4)^k\binom hk\pmod p$,
the sum is $(1+4t)^h$ modulo $p$. Euler's criterion, including
the case $p\mid1+4t$, gives the stated character.
\end{proof}

At $A=42$, $B=37$, $t=-16$, the first sum in
\eqref{eq:general-signed-transfer} is
\[
 K_h=\sum_{k=0}^h(42k+16)b_k^3.
\]
For every odd $p\ne3$, quadratic reciprocity gives
$(-63/p)=(-7/p)=(p/7)$, with both sides zero at $p=7$.
At $p=3$, the replacement of $(-63/p)$ by $(p/7)$ is still
valid in \eqref{eq:general-signed-transfer}, since its multiplier
$21p$ is divisible by $p^2$. Therefore
\eqref{eq:sun57-target} is reduced, for every odd prime, to
$K_h\equiv16p\pmod{p^2}$.

\subsection{A finite triangular WZ certificate}

For $n\ge0$ and $0\le k\le n$, put
\[
 H(n,k)=(-1)^k
 \frac{(2n)!^2(2n-k)!}{n!^3(n-k)!^2(n+k)!}
 =b_n^3\frac{(-n)_k^2}{(n+1)_k(-2n)_k}.
\]
At $(n,k)=(0,0)$ use the factorial expression, and set $H=0$
outside the triangular range. Define
\[
 D(n,k)=\frac{4(21n^2-14nk+k^2+8n-4k)}{2n-k}H(n,k),
 \qquad
 E(n,k)=\frac{16(2n+1)^2}{n+k+1}H(n,k),
\]
with $D(0,0)=16$, and with $D,E$ zero outside the triangle.
The certificate is
\begin{equation}\label{eq:sun57-wz}
 D(n,k+1)-D(n,k)=E(n-1,k)-E(n,k).
\end{equation}
For $0\le k<n$, it follows by substituting the ratios
\[
 \frac{H(n,k+1)}{H(n,k)}
 =-\frac{(n-k)^2}{(n+k+1)(2n-k)},
\]
\[
 \frac{H(n-1,k)}{H(n,k)}
 =\frac{n(n-k)^2(n+k)}
        {4(2n-1)^2(2n-k)(2n-k-1)}
\]
and clearing denominators in the resulting rational identity.
At $k=n\ge1$, the certificate reduces to
\[
 D(n,n)=E(n,n)=16(2n+1)H(n,n).
\]
At $(0,0)$ it is immediate from the definitions. Beyond the
triangle both sides vanish. Thus every boundary is included.

Summing \eqref{eq:sun57-wz} over $0\le n,k\le h$, using
$D(n,h+1)=0$ and $E(-1,k)=0$, gives the exact finite identity
\begin{align}
 \sum_{n=0}^h(42n+16)b_n^3
 &=\sum_{k=0}^h E(h,k)\notag\\
 &=16(2h+1)^2b_h^3
   \sum_{k=0}^h
   \frac{(-h)_k^2}{(h+1)_k(-2h)_k(h+k+1)}.
 \label{eq:sun57-triangular-sum}
\end{align}
The pair is motivated by reflection of Guillera's second WZ pair
for Identity~2 in \cite{Guillera2008}, associated with the
$(42n+5)/64^n$ Ramanujan formula. The displayed certificate has
been proved algebraically, so no analytic continuation is used.

Now take $h=(p-1)/2$. For $k<h$, every displayed denominator in
\eqref{eq:sun57-triangular-sum} is a $p$-adic unit, so the
corresponding $E(h,k)$ lies in $p^2\mathbb Z_{(p)}$. At $k=h$,
\[
 \frac{(-h)_h^2}{(h+1)_h(-2h)_h}
 =\frac{(-1)^h}{b_h^2},
 \qquad
 E(h,h)=16p(-1)^h b_h\equiv16p\pmod{p^2},
\]
because $b_h=\binom{p-1}h\equiv(-1)^h\pmod p$.
Thus $K_h\equiv16p\pmod{p^2}$ for every odd prime.

\subsection{Completion and the prime two}

Combining the signed transfer with the last congruence proves
\eqref{eq:sun57-target} for all odd primes, including $3$ and $7$.
For $p=2$, the original two-term sum is
$37-79=-42\equiv2\pmod4$, while the proposed right side is
$2(21(2/7)+16)=74\equiv2\pmod4$. Hence every prime is covered.

This argument concerns only the prime-square assertion. It does
not, on its own, prove the all-integer divisibility assertion or
the unweighted assertion in Conjecture~5.7(i).


\section{Scope, prior results, and verification}

The proof establishes both assertions of the unchanged Conjecture
5.7(ii). The new all-index argument developed here has two separate
arithmetic tasks: proving the bases \(p^e\mid Y(p^e)\), and proving
that those bases imply \(p^e\mid Y(p^em)\) for every odd cofactor
\(m\). The exact digit kernel and its two neighboring coefficient
states are essential for the latter task. Proving only the first
task would not establish the original all-index conjecture.

The formal coordinate change by itself does not prove the result.
It converts a two-variable integrality problem into a precise
divisibility condition. He's partial-sum theorem pays exactly for
the denominators in the reflected-product differences; the digit
operator then preserves that divisibility through every carry.
The even diagonals are dealt with at all prime powers, including
the prime two, by the separate termwise lemma.

For the prime-square assertion, the finite triangular identity
supplies an independent proof of the required known cubic
congruence. Its earlier occurrence in \cite{SunEuler2011} is
explicitly acknowledged. Sun's polynomial transformation
\cite[Theorem 1.1(i)]{SunApery2012} and Guo--Zeng's weighted
Ap\'ery-polynomial congruence \cite[Theorem 1.1(ii)]{GuoZeng2012}
also yield the weighted transfer after an index shift and
coefficientwise differentiation. Consequently this part should
not be used as evidence for a previously unknown general method.

A targeted literature check of the original \(W_n(1,-16)\)
all-index assertion, its weights, and OEIS A179537 did not locate
an earlier proof of precisely \eqref{eq:main-allN}. Such a search
cannot certify absence from the literature. This manuscript makes
no claim of historical firstness, and does not claim that every
novelty or priority condition has been settled. It also does not
claim a solution of the remaining ordinary-prime branches of
Sun's separate Conjecture 5.8(i).

The proof is symbolic. Exact computations were used only to
check transcription and boundary formulas: the original and
diagonal divisibilities for \(1\le N\le300\); the termwise
even-index lemma for \(0\le i,l\le60\); the exact reflection for
odd \(N\le39\); the two-state identities for
\(p=3,5,7,11\) and odd cofactors \(1,3,5,7\), using three
polynomial state pairs; and independent bivariate diagonal
extraction through \(N=44\). The prime-square check covered
every prime through \(131\). These finite checks are not a
substitute for the unrestricted proofs above.

\appendix
\section{A small reproducible check of the two original statements}

The following Python code uses exact integers and no third-party
packages. It checks a bounded range only. The all-quantifier
arguments are contained in the body of the paper.
\begin{Verbatim}[fontsize=\small]
from math import comb, isqrt

def is_prime(p):
    return p >= 2 and all(p % d for d in range(2, isqrt(p) + 1))

def chi7(p):
    r = p % 7
    return 0 if r == 0 else (1 if r in (1, 2, 4) else -1)

def W(n):
    return sum((-16)**k * comb(n, 2*k)**2 * comb(2*k, k)**2
               for k in range(n//2 + 1))

total = 0
for N in range(1, 301):
    n = N - 1
    total += (42*n + 37) * (-1)**n * W(n)
    assert total % N == 0
    if is_prime(N):
        assert (total - N*(21*chi7(N) + 16)) % (N*N) == 0
print("PASS: original divisibility and prime-square checks")
\end{Verbatim}

\begin{thebibliography}{9}

\bibitem{Sun2010v11}
Z.-W. Sun,
\emph{On Ap\'ery numbers and generalized central trinomial coefficients},
arXiv:1006.2776v11, 19 August 2010,
Conjecture 5.7(ii), p.~32.
\url{https://arxiv.org/abs/1006.2776v11}.

\bibitem{He2017}
B. He,
On the divisibility properties concerning sums of binomial coefficients,
\emph{Ramanujan J.} \textbf{43} (2017), 313--326;
published online 29 March 2016.
\url{https://doi.org/10.1007/s11139-016-9780-6}.
The required \((42k+5)\)-divisibility is displayed explicitly in the
publisher's abstract; no unverified theorem number is assigned here.

\bibitem{SunEuler2011}
Z.-W. Sun,
Super congruences and Euler numbers,
\emph{Sci. China Math.} \textbf{54} (2011), 2509--2535,
Theorem 1.3, equation (1.15).
\url{https://arxiv.org/abs/1001.4453}.

\bibitem{SunApery2012}
Z.-W. Sun,
On sums of Ap\'ery polynomials and related congruences,
\emph{J. Number Theory} \textbf{132} (2012), 2673--2699.
\url{https://doi.org/10.1016/j.jnt.2012.05.014}.
\url{https://arxiv.org/abs/1101.1946}.

\bibitem{GuoZeng2012}
V. J. W. Guo and J. Zeng,
Proof of some conjectures of Z.-W. Sun on congruences for Ap\'ery polynomials,
\emph{J. Number Theory} \textbf{132} (2012), 1731--1740.
\url{https://doi.org/10.1016/j.jnt.2012.02.004}.
\url{https://arxiv.org/abs/1101.0983}.

\bibitem{Guillera2008}
J. Guillera,
Hypergeometric identities for 10 extended Ramanujan-type series,
\emph{Ramanujan J.} \textbf{15} (2008), 219--234.
\url{https://arxiv.org/abs/1104.0396}.

\end{thebibliography}
\end{document}
